% Generated by generate_artifacts.py; do not edit.
\newcommand{\ExperimentGeneratedAt}{2026-08-24T03:13:55.088107+00:00}

\begin{table*}[t]
\centering
\scriptsize
\setlength{\tabcolsep}{2.4pt}
\caption{Single-workload strategy comparison. Utilization is logical busy time divided by logical elapsed time.}
\label{tab:strategy-baseline}
\begin{tabular}{lrrrrrrrrrrr}
\toprule
Scheduler & $\lambda$ & Loops & Req./sched. & Sched. ms & Queue & Exec. & E2E & P99 & Cache & Util. & QPS \\
\midrule
random & -- & -- & 6.5 & 0.014 & 75767.26 & 353.10 & 76140.78 & 148277.11 & 13.0 & 99.7 & 113.8 \\
max\_cache\_hit & -- & -- & 6.5 & 0.094 & 87278.33 & 360.32 & 87659.14 & 537843.90 & 43.7 & 99.8 & 56.3 \\
cost & 0.5 & 50.0 & 6.5 & 6.150 & 2.78 & 333.82 & 363.28 & 426.20 & 31.2 & 71.6 & 178.1 \\
\bottomrule
\end{tabular}
\par\vspace{3pt}
\begin{minipage}{0.98\textwidth}
\footnotesize The table compares scheduler overhead and service outcomes under the same single-workload trace to test whether joint cost-aware routing avoids overload. CostScheduler preserves the offered 178.1 QPS with low queueing, whereas random and cache-only routing build long queues despite lower overhead or higher cache hit rate.
\end{minipage}
\end{table*}

\begin{table*}[t]
\centering
\scriptsize
\setlength{\tabcolsep}{2.4pt}
\caption{Trace composition and per-source outcomes in the unified CostScheduler run at lambda zero.}
\label{tab:workloads}
\begin{tabular}{lrrrrrrrr}
\toprule
Trace & Requests & Share & Mean & P50 & P90 & P99 & Cache & Batch \\
\midrule
agent\_biz\_4091 & 51815 & 57.4 & 6465.0 & 2746 & 21270 & 30828 & 17.2 & 2.22 \\
agent\_biz\_69eb & 19819 & 22.0 & 65014.4 & 59188 & 103802 & 124609 & 7.1 & 1.02 \\
agent\_biz\_aae5 & 4504 & 5.0 & 92667.1 & 70789 & 150959 & 409817 & 6.2 & 1.02 \\
agent\_biz\_f233 & 10299 & 11.4 & 8277.0 & 4272 & 23576 & 30874 & 14.2 & 2.96 \\
agent\_biz\_ddec & 3854 & 4.3 & 234290.2 & 186141 & 395937 & 786440 & 0.3 & 1.01 \\
\bottomrule
\end{tabular}
\par\vspace{3pt}
\begin{minipage}{0.98\textwidth}
\footnotesize The table characterizes the five source traces mixed into one pool so aggregate results can be interpreted against workload skew. Mean prompt length spans 6.5k to 234k tokens, and the longest source has both the lowest cache hit rate and almost no batching.
\end{minipage}
\end{table*}

\begin{table*}[t]
\centering
\scriptsize
\setlength{\tabcolsep}{2.4pt}
\caption{Objective ablation on identical schedule contexts. All loop counts are measured rather than configured values.}
\label{tab:objective-ablation}
\begin{tabular}{llrrrrrrr}
\toprule
Weight & $\lambda$ & Loops & Queue & Exec. & E2E & P99 & Cache & Batch \\
\midrule
linear & 0.0 & 2 & 3.23 & 331.92 & 355.89 & 567.09 & 22.6 & 1.77 \\
linear & 0.0 & 50 & 24.14 & 325.35 & 376.55 & 619.55 & 16.8 & 1.50 \\
$\log(1+128x)/128$ & 0.0 & 2 & 0.86 & 329.90 & 351.53 & 462.31 & 27.0 & 2.43 \\
$\log(1+128x)/128$ & 0.0 & 5 & 0.35 & 327.60 & 349.01 & 421.02 & 25.5 & 2.32 \\
$\log(1+128x)/128$ & 0.0 & 10 & 0.26 & 327.20 & 349.28 & 415.67 & 24.7 & 2.29 \\
$\log(1+128x)/128$ & 0.0 & 20 & 0.16 & 326.77 & 350.08 & 401.25 & 24.3 & 2.27 \\
$\log(1+128x)/128$ & 0.0 & 50 & 0.25 & 326.59 & 353.83 & 398.65 & 22.8 & 2.22 \\
linear & 0.1 & 50 & 23.82 & 324.73 & 375.41 & 600.97 & 17.8 & 1.49 \\
$\log(1+128x)/128$ & 0.1 & 50 & 0.21 & 326.27 & 353.25 & 392.73 & 23.1 & 2.21 \\
linear & 0.5 & 50 & 60.21 & 323.96 & 411.11 & 711.75 & 19.6 & 1.48 \\
$\log(1+128x)/128$ & 0.5 & 50 & 0.97 & 325.88 & 353.79 & 409.52 & 23.5 & 2.13 \\
\bottomrule
\end{tabular}
\par\vspace{3pt}
\begin{minipage}{0.98\textwidth}
\footnotesize The table compares linear and logarithmic latency weights on identical scheduling contexts to test the dispersion correction. The logarithmic weight prevents queue growth at deeper optimization; five loops gives the lowest E2E at lambda zero, while further loops mainly improve the projected objective.
\end{minipage}
\end{table*}

\begin{table*}[t]
\centering
\scriptsize
\setlength{\tabcolsep}{2.4pt}
\caption{Single-workload CostScheduler sweep.}
\label{tab:single-lambda}
\begin{tabular}{lrrrrrrrrrrr}
\toprule
Run & $\lambda$ & Loops & Req./sched. & Sched. ms & Queue & Exec. & E2E & P99 & Cache & Util. & QPS \\
\midrule
l0 & 0.0 & 50.0 & 6.5 & 6.290 & 2.51 & 333.29 & 362.70 & 424.11 & 29.1 & 70.9 & 178.0 \\
l0\_25 & 0.25 & 50.0 & 6.5 & 6.162 & 2.53 & 333.48 & 362.69 & 424.33 & 30.3 & 71.0 & 178.1 \\
l0\_5 & 0.5 & 50.0 & 6.5 & 6.150 & 2.78 & 333.82 & 363.28 & 426.20 & 31.2 & 71.6 & 178.1 \\
l0\_75 & 0.75 & 50.0 & 6.5 & 6.361 & 3.17 & 334.21 & 364.34 & 427.79 & 32.6 & 72.4 & 178.1 \\
l0\_8 & 0.8 & 50.0 & 6.5 & 6.311 & 3.53 & 334.31 & 364.76 & 430.72 & 32.6 & 72.8 & 178.0 \\
l0\_85 & 0.85 & 50.0 & 6.5 & 6.344 & 4.41 & 334.43 & 365.84 & 439.29 & 33.2 & 73.2 & 178.1 \\
l0\_9 & 0.9 & 50.0 & 6.5 & 6.307 & 7.32 & 334.54 & 368.79 & 487.39 & 33.0 & 73.9 & 178.1 \\
l0\_95 & 0.95 & 50.0 & 6.5 & 6.453 & 16.82 & 335.19 & 379.10 & 731.98 & 34.0 & 75.6 & 178.1 \\
l0\_975 & 0.975 & 50.0 & 6.5 & 6.362 & 21.89 & 335.97 & 384.86 & 832.95 & 34.8 & 76.9 & 178.0 \\
l0\_99 & 0.99 & 50.0 & 6.5 & 6.271 & 33.88 & 338.28 & 399.08 & 933.97 & 37.4 & 80.7 & 177.9 \\
l0\_995 & 0.995 & 50.0 & 6.5 & 6.375 & 61.97 & 340.68 & 429.68 & 1107.03 & 40.0 & 84.6 & 178.1 \\
l0\_999 & 0.999 & 50.0 & 6.5 & 6.309 & 350.59 & 347.37 & 724.85 & 2018.26 & 48.7 & 96.1 & 178.1 \\
l0\_9995 & 0.9995 & 50.0 & 6.5 & 6.347 & 798.58 & 349.20 & 1174.74 & 3093.01 & 52.0 & 98.9 & 177.8 \\
l0\_9997 & 0.9997 & 50.0 & 6.5 & 6.419 & 1331.98 & 349.75 & 1708.78 & 4250.17 & 53.2 & 99.7 & 177.8 \\
l0\_9998 & 0.9998 & 50.0 & 6.5 & 6.394 & 2066.55 & 349.95 & 2443.54 & 5670.13 & 53.5 & 99.7 & 176.9 \\
l0\_9999 & 0.9999 & 50.0 & 6.5 & 6.360 & 4084.39 & 350.14 & 4461.56 & 8755.40 & 53.0 & 99.7 & 175.7 \\
l0\_99995 & 0.99995 & 50.0 & 6.5 & 6.379 & 7414.33 & 349.98 & 7791.37 & 12267.28 & 50.4 & 99.7 & 174.7 \\
l1 & 1.0 & 50.0 & 6.5 & 6.378 & 79727.35 & 360.96 & 80115.35 & 293350.62 & 43.5 & 99.4 & 79.5 \\
\bottomrule
\end{tabular}
\par\vspace{3pt}
\begin{minipage}{0.98\textwidth}
\footnotesize The table sweeps lambda on one workload to show how execution efficiency trades against online queueing. Low lambda gives the best latency region; values near one sharply increase queues and eventually reduce throughput, so no single setting dominates all objectives.
\end{minipage}
\end{table*}

\begin{table*}[t]
\centering
\scriptsize
\setlength{\tabcolsep}{2.4pt}
\caption{Independent pools and one unified mixed-workload pool at equal aggregate capacity.}
\label{tab:pool-baseline}
\begin{tabular}{lrrrrrrrrrrr}
\toprule
Deployment & $\lambda$ & Loops & Req./sched. & Sched. ms & Queue & Exec. & E2E & P99 & Cache & Util. & QPS \\
\midrule
independent random & -- & -- & 3.1 & 0.033 & 63012.94 & 334.92 & 63367.56 & 211717.76 & 1.2 & 45.0 & 165.8 \\
independent max cache hit & -- & -- & 3.1 & 0.059 & 69744.32 & 338.29 & 70102.37 & 317992.14 & 7.4 & 44.2 & 121.9 \\
unified random & -- & -- & 10.2 & 0.015 & 284.94 & 335.00 & 639.58 & 2088.41 & 0.3 & 63.6 & 297.1 \\
unified max cache hit & -- & -- & 10.2 & 0.305 & 8944.16 & 339.86 & 9304.00 & 102416.94 & 9.6 & 60.8 & 197.2 \\
independent cost & 0.0 & 38.0 & 3.1 & 13.266 & 20.79 & 321.34 & 376.63 & 835.95 & 4.5 & 40.8 & 298.4 \\
unified cost & 0.0 & 50.0 & 10.2 & 7.516 & 0.42 & 320.13 & 347.85 & 412.09 & 6.3 & 39.9 & 298.9 \\
\bottomrule
\end{tabular}
\par\vspace{3pt}
\begin{minipage}{0.98\textwidth}
\footnotesize The table holds aggregate capacity fixed while comparing policies and independent versus unified placement, isolating the benefit of shared scheduling. The unified CostScheduler run changes mean E2E from 376.63 ms to 347.85 ms relative to independent CostScheduler without reducing completion throughput.
\end{minipage}
\end{table*}

\begin{table*}[t]
\centering
\scriptsize
\setlength{\tabcolsep}{2.4pt}
\caption{Decision-time total-latency prediction versus posterior observation.}
\label{tab:canonical-prediction}
\begin{tabular}{lrrrr}
\toprule
Case & Samples & Pred. ms & Actual ms & MAPE \\
\midrule
Single, $\lambda=0.5$ & 25656 & 321.21 & 337.73 & 9.34\% \\
Independent, $\lambda=0$ & 58289 & 693.33 & 336.68 & 129.82\% \\
Unified, $\lambda=0$ & 58190 & 673.45 & 315.87 & 128.28\% \\
Unified cap 2, $\lambda=0$ & 51614 & 739.64 & 316.82 & 151.38\% \\
Unified, $\lambda=1$ & 75500 & 2919.25 & 1702.10 & 106.76\% \\
\bottomrule
\end{tabular}
\par\vspace{3pt}
\begin{minipage}{0.98\textwidth}
\footnotesize Prediction is the value available before routing; actual is the posterior queue-plus-execution observation for the matched batch. All unmatched-batch counts are zero. The canonical configuration uses a 100,000-sample fit batch with interval zero, so these runs do not trigger online fitting and evaluate the default static model. The high heterogeneous-workload MAPE is therefore a model-accuracy limitation, not evidence of precise steady-state prediction.
\end{minipage}
\end{table*}

\begin{table*}[t]
\centering
\scriptsize
\setlength{\tabcolsep}{2.4pt}
\caption{CostScheduler sweep with five independent pools.}
\label{tab:independent-lambda}
\begin{tabular}{lrrrrrrrrrrr}
\toprule
Run & $\lambda$ & Loops & Req./sched. & Sched. ms & Queue & Exec. & E2E & P99 & Cache & Util. & QPS \\
\midrule
l0 & 0.0 & 38.0 & 3.1 & 13.266 & 20.79 & 321.34 & 376.63 & 835.95 & 4.5 & 40.8 & 298.4 \\
l0\_1 & 0.1 & 50.0 & 3.1 & 13.177 & 24.04 & 321.39 & 379.78 & 916.38 & 4.4 & 41.5 & 298.5 \\
l0\_25 & 0.25 & 50.0 & 3.1 & 13.224 & 22.85 & 321.44 & 378.68 & 910.74 & 4.4 & 40.9 & 298.5 \\
l0\_5 & 0.5 & 50.0 & 3.1 & 13.387 & 32.32 & 321.60 & 388.47 & 1097.67 & 4.9 & 41.9 & 298.5 \\
l0\_75 & 0.75 & 50.0 & 3.1 & 13.398 & 39.27 & 321.68 & 395.57 & 1270.50 & 5.4 & 41.5 & 298.5 \\
l0\_9 & 0.9 & 50.0 & 3.1 & 13.312 & 53.78 & 321.91 & 410.19 & 1899.53 & 5.3 & 41.8 & 298.5 \\
l0\_95 & 0.95 & 50.0 & 3.1 & 13.297 & 69.49 & 322.15 & 426.15 & 2252.14 & 4.9 & 42.8 & 298.4 \\
l0\_99 & 0.99 & 50.0 & 3.1 & 13.366 & 133.67 & 323.95 & 492.21 & 3333.22 & 4.9 & 43.5 & 298.5 \\
l0\_999 & 0.999 & 50.0 & 3.1 & 13.367 & 1001.27 & 327.14 & 1362.95 & 9910.39 & 6.4 & 46.5 & 295.2 \\
l1 & 1.0 & 50.0 & 3.1 & 13.274 & 113941.34 & 338.68 & 114314.56 & 609166.39 & 6.7 & 41.4 & 88.0 \\
\bottomrule
\end{tabular}
\par\vspace{3pt}
\begin{minipage}{0.98\textwidth}
\footnotesize The table sweeps lambda while each workload retains its own Agent pool, providing a control for the shared-pool experiment. Latency is broadly stable through the middle range, but degrades at 0.99 and collapses at the queue-blind lambda-one boundary.
\end{minipage}
\end{table*}

\begin{table*}[t]
\centering
\scriptsize
\setlength{\tabcolsep}{2.4pt}
\caption{CostScheduler sweep after combining five workloads and their nodes.}
\label{tab:unified-lambda}
\begin{tabular}{lrrrrrrrrrrr}
\toprule
Run & $\lambda$ & Loops & Req./sched. & Sched. ms & Queue & Exec. & E2E & P99 & Cache & Util. & QPS \\
\midrule
l0 & 0.0 & 50.0 & 10.2 & 7.516 & 0.42 & 320.13 & 347.85 & 412.09 & 6.3 & 39.9 & 298.9 \\
l0\_1 & 0.1 & 50.0 & 10.3 & 7.562 & 0.41 & 319.95 & 347.72 & 411.79 & 6.2 & 39.8 & 298.9 \\
l0\_25 & 0.25 & 50.0 & 10.2 & 7.458 & 0.53 & 319.87 & 347.60 & 411.95 & 6.3 & 39.8 & 299.0 \\
l0\_5 & 0.5 & 50.0 & 10.2 & 7.578 & 0.63 & 319.46 & 347.49 & 410.58 & 5.9 & 39.7 & 299.0 \\
l0\_75 & 0.75 & 50.0 & 10.2 & 7.574 & 0.87 & 319.19 & 347.45 & 411.21 & 5.7 & 39.8 & 299.0 \\
l0\_9 & 0.9 & 50.0 & 10.2 & 7.498 & 1.01 & 319.10 & 347.38 & 411.87 & 5.8 & 39.6 & 299.0 \\
l0\_95 & 0.95 & 50.0 & 10.2 & 7.581 & 1.41 & 319.01 & 347.77 & 412.34 & 5.8 & 39.5 & 298.9 \\
l0\_99 & 0.99 & 50.0 & 10.2 & 7.570 & 4.53 & 319.45 & 351.36 & 429.38 & 5.8 & 39.9 & 298.9 \\
l0\_999 & 0.999 & 50.0 & 10.2 & 7.524 & 22.71 & 324.74 & 374.78 & 619.23 & 6.3 & 44.8 & 299.0 \\
l1 & 1.0 & 50.0 & 10.2 & 7.550 & 1386.65 & 332.16 & 1746.14 & 33055.59 & 8.3 & 54.3 & 217.9 \\
\bottomrule
\end{tabular}
\par\vspace{3pt}
\begin{minipage}{0.98\textwidth}
\footnotesize The table applies the same sweep after all workloads share one Agent pool, testing whether the online operating region survives mixing. Low-lambda settings remain close on this trace, while near-one settings build large queues and lose throughput despite lower execution time.
\end{minipage}
\end{table*}

\begin{table*}[t]
\centering
\scriptsize
\setlength{\tabcolsep}{2.4pt}
\caption{Effect of the Master collection window.}
\label{tab:window-sweep}
\begin{tabular}{lrrrrrrrrrrr}
\toprule
Window & $\lambda$ & Loops & Req./sched. & Sched. ms & Queue & Exec. & E2E & P99 & Cache & Util. & QPS \\
\midrule
30 requests / 30 ms & 0.0 & 200.0 & 6.5 & 23.387 & 1.28 & 332.90 & 378.50 & 436.64 & 28.2 & 70.0 & 178.1 \\
50 requests / 50 ms & 0.0 & 200.0 & 10.1 & 23.995 & 0.66 & 327.99 & 383.78 & 447.12 & 31.8 & 63.5 & 178.0 \\
75 requests / 75 ms & 0.0 & 200.0 & 14.6 & 24.436 & 0.50 & 324.06 & 392.82 & 467.44 & 34.2 & 58.9 & 178.3 \\
100 requests / 100 ms & 0.0 & 200.0 & 19.1 & 25.321 & 0.27 & 321.45 & 403.27 & 486.91 & 35.2 & 55.9 & 178.1 \\
\bottomrule
\end{tabular}
\par\vspace{3pt}
\begin{minipage}{0.98\textwidth}
\footnotesize The table varies the request-count and wait-time limits together to expose the batching-versus-delay tradeoff. From the smallest to the largest window, mean batch size rises from 2.136 to 2.648 and card-time falls, but mean E2E rises from 378.50 to 403.27 ms and P99 also increases. The largest window is therefore not uniformly best.
\end{minipage}
\end{table*}

\begin{table*}[t]
\centering
\scriptsize
\setlength{\tabcolsep}{2.4pt}
\caption{Optimization-loop cap under the shared mixed workload.}
\label{tab:loop-sweep}
\begin{tabular}{lrrrrrrrrrrr}
\toprule
Cap & $\lambda$ & Loops & Req./sched. & Sched. ms & Queue & Exec. & E2E & P99 & Cache & Util. & QPS \\
\midrule
cap 50 & 0.0 & 50.0 & 10.2 & 7.516 & 0.42 & 320.13 & 347.85 & 412.09 & 6.3 & 39.9 & 298.9 \\
cap 2 & 0.0 & 2.0 & 10.2 & 0.565 & 0.62 & 318.26 & 339.19 & 408.03 & 6.7 & 35.4 & 298.9 \\
cap 5 & 0.0 & 5.0 & 10.2 & 1.015 & 0.52 & 319.76 & 340.99 & 405.98 & 6.3 & 38.7 & 299.0 \\
cap 10 & 0.0 & 10.0 & 10.2 & 1.764 & 0.45 & 319.96 & 341.91 & 406.47 & 6.3 & 39.4 & 299.0 \\
cap 20 & 0.0 & 20.0 & 10.2 & 3.156 & 0.48 & 320.11 & 343.35 & 407.97 & 6.5 & 39.7 & 299.1 \\
\bottomrule
\end{tabular}
\par\vspace{3pt}
\begin{minipage}{0.98\textwidth}
\footnotesize The table varies only the optimization-loop cap on the shared workload to measure route quality against scheduler cost. In this trace, raising the cap from two to fifty changes mean E2E from 339.19 to 347.85 ms and mean scheduling time from 0.565 to 7.516 ms. Deeper optimization does not improve the service outcome in this run.
\end{minipage}
\end{table*}

\begin{table*}[t]
\centering
\scriptsize
\setlength{\tabcolsep}{2.4pt}
\caption{Distribution of one deduplicated CostScheduler objective value per scheduling round. Losses are comparable only within one lambda and workload configuration.}
\label{tab:loss-distribution}
\begin{tabular}{llrrrrrr}
\toprule
Deployment & $\lambda$ & Rounds & Mean & P50 & P90 & P99 & Max \\
\midrule
Single & 0.0 & 8231 & 9.2389 & 7.0968 & 18.986 & 42.979 & 115.91 \\
Single & 0.25 & 8235 & 6.9192 & 5.3185 & 14.045 & 32.477 & 79.365 \\
Single & 0.5 & 8235 & 4.651 & 3.5672 & 9.464 & 21.744 & 136 \\
Single & 0.75 & 8241 & 2.3761 & 1.8088 & 4.7553 & 11.34 & 99.045 \\
Single & 0.8 & 8242 & 1.9019 & 1.4632 & 3.8401 & 9.2009 & 33.116 \\
Single & 0.85 & 8239 & 1.4563 & 1.1188 & 2.9158 & 6.996 & 85.011 \\
Single & 0.9 & 8233 & 0.98887 & 0.76002 & 2.0132 & 4.9918 & 11.817 \\
Single & 0.95 & 8248 & 0.54203 & 0.41441 & 1.0832 & 2.7039 & 8.0195 \\
Single & 0.975 & 8224 & 0.31535 & 0.23546 & 0.59802 & 1.4031 & 65.996 \\
Single & 0.99 & 8225 & 0.17564 & 0.12714 & 0.31941 & 0.75442 & 63.71 \\
Single & 0.995 & 8248 & 0.12945 & 0.090117 & 0.21641 & 0.56638 & 62.894 \\
Single & 0.999 & 8234 & 0.08614 & 0.060106 & 0.15019 & 0.47782 & 4.8413 \\
Single & 0.9995 & 8243 & 0.091259 & 0.057409 & 0.14546 & 0.46016 & 62.48 \\
Single & 0.9997 & 8225 & 0.083156 & 0.056712 & 0.14146 & 0.43972 & 4.8382 \\
Single & 0.9998 & 8241 & 0.082028 & 0.056259 & 0.14099 & 0.45694 & 4.8413 \\
Single & 0.9999 & 8236 & 0.084071 & 0.056269 & 0.13719 & 0.46873 & 12.479 \\
Single & 0.99995 & 8239 & 0.084143 & 0.055123 & 0.13438 & 0.46629 & 29.219 \\
Single & 1.0 & 8234 & 0.071983 & 0.044818 & 0.12214 & 0.44984 & 6.2719 \\
\midrule
Independent & 0.0 & 29214 & 5.4294 & 0.47821 & 12.596 & 53.464 & 678.17 \\
Independent & 0.1 & 29196 & 5.0494 & 0.43257 & 11.393 & 48.733 & 710.16 \\
Independent & 0.25 & 29179 & 4.1752 & 0.36912 & 9.5519 & 40.007 & 476.28 \\
Independent & 0.5 & 29216 & 2.9961 & 0.25557 & 6.4467 & 30.55 & 456.68 \\
Independent & 0.75 & 29233 & 1.6352 & 0.13493 & 3.417 & 16.608 & 257.44 \\
Independent & 0.9 & 29141 & 0.79709 & 0.065544 & 1.4999 & 7.2077 & 194.46 \\
Independent & 0.95 & 29264 & 0.52544 & 0.042373 & 0.85967 & 4.4273 & 141.6 \\
Independent & 0.99 & 29288 & 0.26405 & 0.023791 & 0.30741 & 1.905 & 86.625 \\
Independent & 0.999 & 29227 & 0.20689 & 0.019403 & 0.188 & 1.2854 & 68.673 \\
Independent & 1.0 & 29316 & 0.16816 & 0.016416 & 0.12253 & 0.82582 & 62.088 \\
\midrule
Unified & 0.0 & 8809 & 6.4512 & 5.115 & 13.41 & 27.15 & 70.5 \\
Unified & 0.1 & 8797 & 5.8181 & 4.6097 & 12.126 & 24.366 & 63.424 \\
Unified & 0.25 & 8826 & 4.861 & 3.8474 & 10.176 & 20.304 & 52.879 \\
Unified & 0.5 & 8815 & 3.2368 & 2.5566 & 6.7573 & 13.604 & 35.422 \\
Unified & 0.75 & 8819 & 1.6306 & 1.2878 & 3.4017 & 6.8773 & 17.779 \\
Unified & 0.9 & 8814 & 0.6648 & 0.52578 & 1.3674 & 2.7198 & 7.1935 \\
Unified & 0.95 & 8817 & 0.34291 & 0.27294 & 0.69696 & 1.3793 & 4.3981 \\
Unified & 0.99 & 8831 & 0.08449 & 0.068648 & 0.16141 & 0.3052 & 3.9415 \\
Unified & 0.999 & 8817 & 0.026072 & 0.022063 & 0.038041 & 0.084635 & 2.2573 \\
Unified & 1.0 & 8824 & 0.019636 & 0.016049 & 0.026646 & 0.076014 & 2.2179 \\
\bottomrule
\end{tabular}
\par\vspace{3pt}
\begin{minipage}{0.98\textwidth}
\footnotesize The table reports the distribution of one deduplicated final objective per scheduling round to expose optimizer spread and outliers. Values are useful within an identical workload and lambda, but their decline as lambda changes must not be interpreted as a service-level improvement.
\end{minipage}
\end{table*}
